\section{The Live Dashboard: What Every Metric Means}
\label{sec:dashboard}

Everything in this report --- the magnitude trick, the Welch spectrum, the Freeze
Index, the still-floor gate, the debounce --- exists to drive \emph{one screen}. When
Cadence is on an ankle and the wearer takes a few steps, a clinician (or the wearer
themselves) looks at a single web page and should be able to say, at a glance,
``\emph{the leg is walking calmly},'' ``\emph{the leg is shuffling in place --- that is a
freeze},'' or ``\emph{the leg is at rest}.'' This section is a reading guide for that
screen. It walks panel by panel, says in plain words what each number means, gives the
\emph{typical} values you see when the leg is still, walking, and freezing, and --- most
importantly --- makes the arithmetic that links the panels explicit, so the dashboard
reads as one coherent story rather than seven disconnected dials.

\begin{keyidea}
The dashboard is a \textbf{glass box}, not a black box. Every panel is a quantity you
can compute by hand from the raw trace, and the headline verdict is built from them by a
rule you can recite: \emph{the gate decides if the leg is moving at all; the Freeze
Index decides whether that movement is a freeze; the debounce makes sure the verdict is
not a one-off twitch}. Nothing on the screen is a learned weight you cannot inspect.
\end{keyidea}

\subsection{The screen at a glance}

Figure~\ref{fig:dash-mock} is an annotated mock of the live page. There are seven
panels. One of them --- \emph{Gait Status} --- is the answer; the other six are the
\emph{evidence} that produced it. The trick to reading the dashboard quickly is to know
which panels are inputs and which is the output, and how the inputs flow into the
output. That flow is always the same three steps, drawn as the spine of the mock:
\[
\underbrace{\textsf{Movement Energy}}_{\text{is the leg moving?}}
\;\longrightarrow\;
\underbrace{\textsf{Freeze Index}}_{\text{is it a freeze?}}
\;\longrightarrow\;
\underbrace{\textsf{Debounce}}_{\text{is it stable?}}
\;\longrightarrow\;
\textsf{Gait Status}.
\]

\begin{figure}[htbp]
\centering
\resizebox{0.97\textwidth}{!}{%
\begin{tikzpicture}[
  font=\sffamily,
  panel/.style={draw=rule,line width=0.8pt,rounded corners=3pt,fill=mist,
                minimum width=4.6cm,minimum height=2.05cm,align=center,inner sep=5pt},
  hdr/.style={font=\footnotesize\bfseries,text=navy},
  sub/.style={font=\scriptsize,text=mute},
  big/.style={font=\Large\bfseries},
  call/.style={font=\footnotesize\itshape,text=ink,align=left},
  arr/.style={-{Stealth[length=2.2mm]},line width=0.9pt,color=steel},
  flow/.style={-{Stealth[length=3mm]},line width=1.6pt,color=navy},
]

% ---------- headline verdict (top, spanning) ----------
\node[panel,fill=palefreeze,draw=freeze,line width=1.2pt,
      minimum width=15.4cm,minimum height=1.9cm] (status) at (0,0) {};
\node[hdr,anchor=north west] at ([xshift=6pt,yshift=-5pt]status.north west) {GAIT STATUS};
\node[big,text=freeze,anchor=west] (verdictword) at ([xshift=-6.6cm,yshift=-2pt]status.center) {FREEZE};
\node[circle,fill=freeze,minimum size=11pt,inner sep=0pt]
  at ([xshift=0.55cm]verdictword.east){};
% thin divider
\draw[rule,line width=0.7pt] ([xshift=-1.9cm,yshift=12pt]status.center) -- ([xshift=-1.9cm,yshift=-12pt]status.center);
\node[sub,align=left,anchor=west] at ([xshift=-1.6cm,yshift=-2pt]status.center)
  {\textbf{STILL} \;\textbullet\; \textbf{WALKING} \;\textbullet\; \textbf{FREEZE}\\[2pt]
   the one headline verdict\\[1pt] gate $\to$ FI $\to$ debounce};

% ---------- Freeze Margin meter (second row, left) ----------
\node[panel,below left=1.0cm and 0.15cm of status.south,anchor=north,
      fill=palefreeze,draw=freeze] (margin) {};
\node[hdr,anchor=north west] at ([xshift=6pt,yshift=-5pt]margin.north west){FREEZE MARGIN};
\node[big,text=freeze] at ([yshift=6pt]margin.center){118\,\%};
% mini meter along the panel floor
\draw[draw=mute,fill=white,line width=0.6pt]
  ([xshift=10pt,yshift=6pt]margin.south west) rectangle ([xshift=130pt,yshift=15pt]margin.south west);
\draw[fill=freeze,draw=none]
  ([xshift=10pt,yshift=6pt]margin.south west) rectangle ([xshift=128pt,yshift=15pt]margin.south west);
\draw[dashed,line width=0.7pt,color=ink]
  ([xshift=80pt,yshift=2pt]margin.south west) -- ([xshift=80pt,yshift=19pt]margin.south west);
\node[font=\tiny,text=mute,anchor=south] at ([xshift=80pt,yshift=18pt]margin.south west){trip};

% ---------- Freeze Index (second row, mid) ----------
\node[panel,right=0.3cm of margin] (fi) {};
\node[hdr,anchor=north west] at ([xshift=6pt,yshift=-5pt]fi.north west){FREEZE INDEX};
\node[big,text=ink] at ([yshift=3pt]fi.center){2.14};
\node[sub] at ([yshift=-16pt]fi.center){ratio \;\textbullet\; trip line $1.815$};

% ---------- Movement Energy (second row, right) ----------
\node[panel,right=0.3cm of fi,fill=paleteal,draw=teal] (energy) {};
\node[hdr,anchor=north west,text=teal] at ([xshift=6pt,yshift=-5pt]energy.north west){MOVEMENT ENERGY};
\node[big,text=teal] at ([yshift=3pt]energy.center){high};
\node[sub] at ([yshift=-16pt]energy.center){0.5--8\,Hz \;\textbullet\; gate: \textbf{MOVING}};

% ---------- third row ----------
\node[panel,below=0.9cm of margin.south,anchor=north,
      minimum width=4.6cm] (loco) {};
\node[hdr,anchor=north west] at ([xshift=6pt,yshift=-5pt]loco.north west){LOCOMOTION POWER};
\node[big,text=ink] at ([yshift=2pt]loco.center){low};
\node[sub] at ([yshift=-16pt]loco.center){0.5--3\,Hz \;\textbullet\; walking rhythm};

\node[panel,right=0.3cm of loco,fill=paleamber,draw=amber] (tremor) {};
\node[hdr,anchor=north west,text=amber] at ([xshift=6pt,yshift=-5pt]tremor.north west){TREMOR POWER};
\node[big,text=amber] at ([yshift=2pt]tremor.center){present};
\node[sub] at ([yshift=-16pt]tremor.center){4--6\,Hz \;\textbullet\; rest-tremor band};

\node[panel,right=0.3cm of tremor] (wave) {};
\node[hdr,anchor=north west] at ([xshift=6pt,yshift=-5pt]wave.north west){ACCELEROMETER};
% tiny jittery trace, drawn relative to the panel centre
\begin{scope}[shift={([yshift=2pt]wave.center)}]
  \draw[line width=0.5pt,color=steel] plot[smooth,tension=0.6]
    coordinates {(-1.6,-0.12)(-1.3,0.2)(-1.0,-0.24)(-0.7,0.28)(-0.4,-0.16)(-0.1,0.32)(0.2,-0.28)(0.5,0.24)(0.8,-0.2)(1.1,0.26)(1.45,-0.14)};
\end{scope}
\node[sub] at ([yshift=-18pt]wave.center){raw 3-axis @ 64\,Hz};

% ---------- spine flow arrows ----------
% second-row evidence flows right-to-left: gate (energy) -> FI -> margin
\draw[flow] (energy) -- (fi);
\draw[flow] (fi) -- (margin);
% one clean arrow rises from the margin meter up into the verdict banner
\draw[flow] (margin.north) -- ([yshift=-0.02cm]margin.north|-status.south);

% ---------- callouts ----------
\node[call,text width=3.5cm,anchor=west] (c1) at ([xshift=0.35cm]status.east|-status.center){};
\node[call,text width=3.6cm] at ([xshift=2.0cm,yshift=0.0cm]status.east)
  {\textbf{The verdict.}\\ Red dot + word.\\ Built from the three panels below.};

\node[call,text width=3.4cm] at ([xshift=-2.2cm,yshift=0.0cm]margin.west)
  {\textbf{How close to a freeze?}\\ FI as \% of the trip.\\ $>$100\,\% $\Rightarrow$ over.};

\node[call,text width=3.4cm] at ([xshift=2.25cm,yshift=0.0cm]energy.east)
  {\textbf{The gate.}\\ Below floor $\Rightarrow$ forced STILL,\\ margin shows ``\textemdash/at rest''.};

\node[call,text width=3.4cm] at ([xshift=2.25cm,yshift=0.0cm]wave.east)
  {\textbf{Ground truth.}\\ The shuffle you would\\ see by eye, in numbers.};

\end{tikzpicture}}
\caption{Annotated mock of the live web dashboard. The red \textsc{Gait Status} banner
is the single verdict; the six panels beneath it are the evidence. Heavy navy arrows
trace the decision spine: \emph{Movement Energy} (the gate) feeds the \emph{Freeze
Index}, which feeds the \emph{Freeze Margin} meter and, after debounce, the verdict.
Colours follow the report palette: teal $=$ sensing/energy, amber $=$ tremor, red $=$
the freeze decision.}
\label{fig:dash-mock}
\end{figure}

\begin{intuition}
Think of the panels as a courtroom. \emph{Movement Energy} is the bailiff at the door:
if the leg is barely moving, the case is dismissed before it starts and the verdict is
\textbf{STILL}. \emph{Freeze Index} is the prosecutor: given that the leg \emph{is}
moving, it argues whether the movement is the fast 3--8\,Hz shuffle of a freeze or the
slower swing of a walk. \emph{Debounce} is the jury foreman, refusing to convict on a
single noisy instant. \emph{Gait Status} reads out the sentence.
\end{intuition}

\subsection{The two quantities everything is built from}

Before the panels make sense, two numbers must be clear, because every other panel is a
slice or a combination of them. Both are computed on the mean-removed acceleration
\emph{magnitude} $a=\sqrt{a_x^2+a_y^2+a_z^2}$ over the current 4-second window, and both
are sums of the Welch power spectrum over a frequency band.

\begin{mathsbox}
\textbf{Band power} is the area under the power spectrum across a band --- literally how
much wiggle energy lives at those frequencies. Cadence reads three bands:
\[
P_{\text{loco}} = P(0.5\text{--}3\,\text{Hz}),\quad
P_{\text{freeze}} = P(3\text{--}8\,\text{Hz}),\quad
P_{\text{tremor}} = P(4\text{--}6\,\text{Hz}).
\]
From these, the two master quantities are
\[
\boxed{\;\FI = \dfrac{P_{\text{freeze}}}{P_{\text{loco}}+10^{-9}}\;}
\qquad\text{and}\qquad
\boxed{\;E_{\text{move}} = P_{\text{loco}} + P_{\text{freeze}}\;}.
\]
$\FI$ is a \emph{ratio} (the freeze band \emph{divided by} the walking band), so it is
scale-invariant: it does not care whether you measure in mg or g. $E_{\text{move}}$ is a
\emph{sum} (the two bands \emph{added}), so it grows with how hard the leg is actually
moving and is \emph{not} scale-invariant. The tiny $10^{-9}$ stops a perfectly still leg
from dividing by zero.
\end{mathsbox}

This single pair of formulas is the Rosetta Stone for the whole dashboard:

\begin{itemize}
\item \textbf{Locomotion Power} \emph{is} $P_{\text{loco}}$ --- the denominator of $\FI$
  and one half of $E_{\text{move}}$.
\item \textbf{Tremor Power} \emph{is} $P_{\text{tremor}}$ --- a narrow sub-band sitting
  \emph{inside} the freeze band.
\item \textbf{Freeze Index} \emph{is} $\FI$ --- freeze band over loco band.
\item \textbf{Movement Energy} \emph{is} $E_{\text{move}}$ --- freeze band plus loco band.
\item \textbf{Freeze Margin} re-expresses $\FI$ as a percentage of the trip line.
\item \textbf{Gait Status} combines the gate (on $E_{\text{move}}$), $\FI$, and debounce.
\end{itemize}

\begin{keyidea}
Same two ingredients, different recipe. \emph{Divide} the freeze band by the walking
band and you get the \textbf{shape} of the motion ($\FI$ --- \emph{is this fast or
slow?}). \emph{Add} them and you get the \textbf{amount} of motion ($E_{\text{move}}$
--- \emph{is the leg moving at all?}). The dashboard needs both: amount to know whether
to bother, shape to know whether it is a freeze.
\end{keyidea}

\subsection{Panel by panel}

\subsubsection*{Gait Status --- the verdict}

This is the only panel a hurried reader needs. It shows one of three words with a
matching NeoPixel colour: \textbf{STILL} (the leg is at rest), \textbf{WALKING} (ordinary
gait), or \textbf{FREEZE} (red --- a freeze-of-gait episode is happening \emph{now}).
The verdict is not read off any single dial; it is the output of the three-step state
machine:
\begin{enumerate}
\item \textbf{Gate.} If $E_{\text{move}}$ is below the auto-calibrated still-floor, the
  leg is not really moving $\Rightarrow$ \textbf{STILL}. Stop here.
\item \textbf{Classify.} Otherwise, if $\FI > 1.815$, the moving leg is shuffling fast
  $\Rightarrow$ candidate \textbf{FREEZE}; else $\Rightarrow$ candidate \textbf{WALKING}.
\item \textbf{Debounce.} A new candidate must persist for \emph{two} consecutive
  2-second control windows before it is committed (assert-after-2, release-after-2).
\end{enumerate}

\begin{intuition}
Why debounce? One heel-strike, one stumble, one bump against a table leg can spike $\FI$
for a single window. Without debounce the banner would flicker FREEZE for two seconds
every time the wearer kicked a kerb. Requiring two windows in a row means a real
verdict is at most $\sim$2\,s late but is never a phantom.
\end{intuition}

\subsubsection*{Freeze Margin --- the glass-box confidence}

This meter is the newest panel and the most useful for ``how worried should I be?''. It
expresses the live $\FI$ as a percentage of the $1.815$ trip line:
\[
\text{margin} = \frac{\FI}{1.815}\times 100\,\%.
\]
At \textbf{100\,\%} the leg is exactly on the trip line; \textbf{above 100\,\%} it has
crossed into freeze territory; well below 100\,\% the gait is comfortably calm. Crucially,
when the gate has forced the state to STILL, the meter does \emph{not} show a misleading
small percentage --- it shows ``\,\textemdash\,/ at rest\,'', because a freeze ratio computed on a
motionless leg is meaningless noise.

\begin{keyidea}
The Freeze Margin replaces the old, blank ``probability'' box. In glass-box (FI) mode
the CNN is \emph{bypassed} entirely, so there was never a neural-network probability to
display --- the box sat empty and confused people. The margin meter fills that slot with
something the FI detector can actually justify: a transparent ``percent of the way to a
freeze'' that you can recompute by hand from $\FI$ and $1.815$.
\end{keyidea}

\subsubsection*{Freeze Index --- the raw ratio}

The same $\FI$ as a bare number, with the trip line drawn at $1.815$. A calm walk sits
around $\FI\approx 0.5$--$1$; a genuine freeze pushes well above $2$, and on real
captures the freeze peaks reach into the tens (the worksheet capture peaked at $21.35$).
The dashboard caps the displayed scale so a single off-the-chart spike does not flatten
the rest of the trace --- the same off-scale behaviour visible in
Figure~\ref{fig:smoke}, where one freeze window's $\FI$ runs clean off the top of the
axis.

\subsubsection*{Tremor Power --- the rest-tremor band}

This panel reports $P_{\text{tremor}}$, the power in the narrow 4--6\,Hz rest-tremor band,
with a \textbf{present / clear} badge. Parkinsonian rest tremor lives here, partly
overlapping the lower freeze band. It is reported as a separate readout (not folded into
the verdict) so a clinician can see tremor activity independently of the freeze decision
--- the leg can tremble without freezing, and vice versa.

\subsubsection*{Movement Energy --- the gate}

This is $E_{\text{move}} = P_{\text{loco}}+P_{\text{freeze}}$, the 0.5--8\,Hz energy,
shown with a \textbf{MOVING / STILL} badge that fires when the energy crosses the
auto-calibrated still-floor. It answers one question only: \emph{is the leg doing
anything at all?} A still leg has near-zero energy; the gate is what stops sensor noise
on a motionless ankle from producing a large, meaningless $\FI$ and a phantom freeze.
Because $E_{\text{move}}$ is a raw sum (not a ratio), it is the one panel whose absolute
value genuinely tracks ``how hard the leg is moving.''

\subsubsection*{Locomotion Power --- the walking rhythm}

This is $P_{\text{loco}}$, the 0.5--3\,Hz power: the energy of the slow, rhythmic leg
swing of ordinary walking. It is high during a steady walk and low at rest. It is also
the \emph{denominator} of $\FI$ --- which is exactly why a freeze stands out: during a
freeze the walking rhythm collapses ($P_{\text{loco}}$ falls) while the fast shuffle
grows ($P_{\text{freeze}}$ rises), so their ratio $\FI$ shoots up.

\subsubsection*{Accelerometer waveform --- the ground truth}

The live raw 3-axis trace at 64\,Hz. This is the unprocessed reality every other panel is
distilled from: calm gait looks like clean periodic swings, a freeze looks like a burst
of fast, low-amplitude trembling ``in place,'' and rest looks nearly flat. A skilled eye
can read a freeze straight off this trace; the other panels exist to do that reading
automatically and continuously.

\begin{table}[htbp]
\centering
\caption{Reading the dashboard. Each panel, what it measures, the one-line intuition,
and the values you typically see when the leg is still, walking, or freezing. ``band''
refers to the Welch power spectrum of the mean-removed magnitude.}
\label{tab:panels}
\small
\renewcommand{\arraystretch}{1.25}
\begin{tabular}{@{}p{2.5cm} p{3.0cm} p{3.4cm} >{\raggedright\arraybackslash}p{5.3cm}@{}}
\toprule
\textbf{Panel} & \textbf{What it shows} & \textbf{Intuition} & \textbf{Still / Walk / Freeze} \\
\midrule
\textbf{Gait Status}
  & STILL / WALKING / FREEZE
  & The verdict: gate $\to\FI\to$ debounce
  & STILL~/ WALKING~/ \textcolor{freeze}{\textbf{FREEZE}} \\
\textbf{Freeze Margin}
  & $\FI$ as \% of the $1.815$ trip
  & How close to a freeze (glass-box confidence)
  & ``\textemdash/at rest'' / $\sim$30--55\,\% / $\ge$100\,\% \\
\textbf{Freeze Index}
  & $\FI = P_{\text{freeze}}/P_{\text{loco}}$
  & Fast shuffle vs.\ slow swing (a ratio)
  & gated / $\approx$0.5--1 / $>$2, up to 20+ \\
\textbf{Tremor Power}
  & 4--6\,Hz band power
  & Rest-tremor activity, reported separately
  & clear / clear / present (often) \\
\textbf{Movement Energy}
  & $P_{\text{loco}}+P_{\text{freeze}}$ (0.5--8\,Hz)
  & How hard the leg moves --- the gate
  & STILL (low) / MOVING (high) / MOVING \\
\textbf{Locomotion Power}
  & 0.5--3\,Hz band power
  & The walking-rhythm energy
  & low / high / collapses \\
\textbf{Accelerometer}
  & raw 3-axis trace @ 64\,Hz
  & The shuffle you would see by eye
  & flat / clean swings / fast tremble \\
\bottomrule
\end{tabular}
\end{table}

\subsection{Reading a real run off the screen}

Figure~\ref{fig:smoke} is a recorded smoketest, and it shows in one picture exactly what
the dashboard animates live. The top trace is the accelerometer magnitude through four
labelled phases; the bottom trace is the Freeze Index against its trip line. Walk the
panels through it:

\begin{itemize}
\item \textbf{Phase 0 (rest).} The magnitude is nearly flat. \emph{Movement Energy} is
  below the still-floor, so the gate fires, \emph{Gait Status} reads \textbf{STILL}, and
  the \emph{Freeze Margin} shows ``\textemdash/at rest'' --- regardless of whatever noisy $\FI$ the
  bottom trace happens to show, because a gated leg is never called a freeze.
\item \textbf{Phase 1 (walking).} The magnitude breaks into large, regular swings.
  \emph{Locomotion Power} (and hence \emph{Movement Energy}) rises, so the gate opens and
  the leg is \textbf{MOVING}. But the energy is dominated by the slow walking rhythm, so
  $\FI$ mostly sits at or below the dashed trip line: \emph{Gait Status} reads
  \textbf{WALKING}.
\item \textbf{Phases 2--3 (freeze).} The magnitude turns into bursts of fast, lower-amplitude
  trembling. Now the freeze band swells while the walking band collapses, so $\FI$ spikes
  hard --- one window even runs off the top of the axis (a freeze peak the figure marks
  as off-scale). The \emph{Freeze Margin} crosses 100\,\%, and after two consecutive
  windows confirm it the debounce commits the verdict: \emph{Gait Status} turns red,
  \textbf{FREEZE}.
\end{itemize}

\begin{figure}[htbp]
\centering
\includegraphics[width=0.92\linewidth]{smoketest_trace}
\caption{A recorded smoketest, exactly the signals the dashboard animates live.
\emph{Top:} accelerometer magnitude $\norm{a}$ in mg across four phases --- flat at rest,
big regular swings while walking, fast low-amplitude bursts during a freeze.
\emph{Bottom:} the Freeze Index against its trip line (dashed). At rest the gate forces
STILL; walking keeps $\FI$ near or below the line; the freeze drives $\FI$ far above it,
with one window's peak running clean off the top of the axis. Reading the two traces
together is exactly what the seven panels do, automatically, every two seconds.}
\label{fig:smoke}
\end{figure}

\begin{takeaway}
To read the leg off the screen in three seconds: glance at \textbf{Movement Energy} (is
it moving?); if yes, glance at the \textbf{Freeze Margin} (over 100\,\% $=$ freeze); the
\textbf{Gait Status} banner has already done both, plus debounce, and shows the answer in
one word and one colour. Everything else on the page is there to let you check that
answer by hand.
\end{takeaway}
